% Осень26!A19; нумерация внутри столбца без пустых строк.
% Источник: лекция 7, §11 (строка 365)
\ticket{17}{Теорема о неявном отображении}
\label{ticket:17}

\materialnote{В исходном конспекте финальная проверка открытости и выбора окрестностей отложена. Доказательство требует сверки.}

\subsection*{Теорема о неявном отображении}

\subsubsection*{Соглашения о размерностях; блочная запись производного оператора}

Работаем в \(\mathbb R^{m+n}\); точку записываем как пару \[(x,y),\qquad x=(x_1,\dots,x_m)\in\mathbb R^m,\quad y=(y_1,\dots,y_n)\in\mathbb R^n\] --- «рядышком подряд» \(m+n\) координат. Пусть \(F\colon O\to\mathbb R^n\) --- отображение из открытого \(O\subset\mathbb R^{m+n}\) с координатными функциями \(F_1,\dots,F_n\). Производный оператор \(F'(x,y)\) --- матрица размера \(n\times(m+n)\); в её столбцах --- сначала производные по иксам, потом по игрекам. Мысленно поделив матрицу на два блока, обозначаем: \[F'(x,y)=\bigl(\,F'_x(x,y)\ \bigl|\ F'_y(x,y)\,\bigr),\qquad
F'_x=\Bigl(\frac{\partial F_i}{\partial x_j}\Bigr)_{\substack{i=1..n\\ j=1..m}},\quad
F'_y=\Bigl(\frac{\partial F_i}{\partial y_j}\Bigr)_{\substack{i=1..n\\ j=1..n}} .\] Блок \(F'_y\) --- \textbf{квадратный} \(n\times n\); о нём можно спрашивать, обратим ли он (равен ли нулю его определитель).

\subsubsection*{Мотивирующий пример: окружность}

\(m=n=1\): в \(\mathbb R^2\) уравнение \[F(x,y):=x^2+y^2-25=0\] задаёт окружность радиуса \(5\) (лектор нарисовал картинку). \textbf{Хорошая точка} \((x_0,y_0)=(4,3)\): лежит на окружности, и если брать иксы, близкие к \(4\), то игреки, близкие к \(3\), вычисляются по уравнению \textbf{однозначно} (выбор ветви есть --- значения и около \(-3\), --- но, зафиксировав точку, фиксируем ветвь); по виду линии (окружность «ровненькая, без зубчиков») зависимость \(y(x)\) выглядит гладкой. Саму формулу \(y(x)\) лектор «вычислять не умеет» («квадратное уравнение в мою голову не загрузилось»), но, \textbf{поверив} в дифференцируемость \(y(x)\), можно подставить \(y=y(x)\) в уравнение и продифференцировать: \[2x+2y\,y'=0\quad\Longrightarrow\quad y'=-\frac{x}{y};\qquad y'(4)=-\frac43\] --- угловой коэффициент касательной в точке \((4,3)\) вычисляется, хотя явной формулы \(y(x)\) нет. \textbf{Плохая точка} \((-5,0)\): правее \(-5\) у уравнения два близких решения, левее --- ни одного; зависимость \(y(x)\) в окрестности не строится. И в формуле \(y'=-x/y\) «сбой»: знаменатель \(\partial F/\partial y=2y\) в точке \((-5,0)\) равен нулю, а в точке \((4,3)\): \(2\cdot3=6\ne0\). \textbf{Общий принцип: условие \(\det F'_y\ne0\) --- критерий «хорошей» точки.} («Если взять уравнение посложнее, никакие квадратные уравнения не спасут --- вы просто не сумеете его решить; теорема же гарантирует, что зависимость есть и гладкая.»)

\subsubsection*{Формулировка теоремы}

\begin{theorem}

\emph{о неявном отображении.} Пусть \(O\subset\mathbb R^{m+n}\) открыто, \(F\colon O\to\mathbb R^n\) --- класса \(C^r\), \(r\in\{1,2,\dots\}\cup\{\infty\}\). Пусть найдена точка \((a,b)\in O\) (\(a\in\mathbb R^m\), \(b\in\mathbb R^n\)), в которой \[F(a,b)=0,\] и блок производного оператора по игрекам невырожден: \[\det F'_y(a,b)\ \ne\ 0.\] Тогда:

\textbf{(1)} существуют открытое множество \(P\subset\mathbb R^m\), \(a\in P\), открытое множество \(Q\subset\mathbb R^n\), \(b\in Q\), и отображение \(\varphi\colon P\to Q\) класса \(C^r\), такие что \[F\bigl(x,\varphi(x)\bigr)=0\qquad\text{для всех }x\in P;\]

\textbf{(2)} \(\varphi\) определено однозначно: для каждого \(x\in P\) точка \(\varphi(x)\) --- \textbf{единственный} \(y\in Q\) с \(F(x,y)=0\) («в целом решений может быть и несколько, но те, что близки к \(b\), задают правило однозначно» --- уравнение \(F=0\) на \(P\times Q\) высекает в точности график \(\varphi\));

\textbf{(3)} справедлива формула для производного оператора неявного отображения: \[\varphi'(x)\ =\ -\Bigl(F'_y\bigl(x,\varphi(x)\bigr)\Bigr)^{-1}\,F'_x\bigl(x,\varphi(x)\bigr)\qquad(x\in P)\] («формула не сказать, что очень ценная, она автоматически получается; главное --- виден характер формулы, и в точке \((a,b)\) она полностью вычислима»).

\(\varphi\) и есть \textbf{неявное отображение}: «игреки выражаются через иксы» без явной формулы. Число уравнений (\(n\)) совпадает с числом игреков; уравнение \(F=0\) --- «система из \(n\) уравнений от \(m+n\) переменных».

\end{theorem}

\subsubsection*{Доказательство формулы (3)}

Продифференцируем тождество \(F(x,\varphi(x))=0\). Отображение \(x\mapsto\bigl(x,\varphi(x)\bigr)\) действует из \(\mathbb R^m\) в \(\mathbb R^{m+n}\), его производный оператор --- блочный «столбец» \(\binom{E}{\varphi'(x)}\) (\(m\) столбцов; сверху единичная матрица \(E=E_m\), снизу \(\varphi'\)). По правилу дифференцирования композиции: \[\bigl(\,F'_x\ \bigl|\ F'_y\,\bigr)\binom{E}{\varphi'(x)}\ =\ F'_x\cdot E+F'_y\cdot\varphi'(x)\ =\ 0,\] где все блоки вычислены в точке \(\bigl(x,\varphi(x)\bigr)\). Переносим первое слагаемое вправо и домножаем слева на \((F'_y)^{-1}\): \[\varphi'(x)=-\bigl(F'_y\bigr)^{-1}F'_x.\qquad\blacksquare_{(3)}\] «Чисто возня с формулами, она не особо полезная; получается автоматически, как только доказан пункт 1.»

\subsubsection*{Доказательство существования (пункт 1)}

\textbf{Шаг 1 (подъём размерности).} Определим \(\Phi\colon O\to\mathbb R^{m+n}\) (то же открытое \(O\subset\mathbb R^{m+n}\), теперь в \(\mathbb R^{m+n}\) же и отображаем) правилом: первые \(m\) координат переписываем как есть, вместо остальных пишем значения \(F\): \[\Phi(x,y):=\bigl(x,\ F(x,y)\bigr).\] Тогда \(\Phi(a,b)=(a,0)\) (\(n\)-мерный ноль во второй части). Производный оператор \(\Phi'\) (размера \((m+n)\times(m+n)\)) в блочной форме: \[\Phi'(x,y)=\begin{pmatrix} E_m & 0\\ F'_x(x,y) & F'_y(x,y)\end{pmatrix}\] --- блочно-треугольная матрица, поэтому \[\det\Phi'(a,b)=\det E_m\cdot\det F'_y(a,b)=\det F'_y(a,b)\ne0.\]

\textbf{Шаг 2 (локальная обратимость).} Теперь \(\Phi\) --- отображение «из пространства в пространство той же размерности» с невырожденным производным оператором. По теореме о локальной обратимости (§10) существует окрестность \(W\) точки \((a,b)\) такая, что \(\Phi\big|_W\) --- диффеоморфизм на свой образ \(V:=\Phi(W)\) (открытое множество). Заузим окрестность: диффеоморфизм останется диффеоморфизмом при сужении, поэтому считаем \[W=P_1\times Q,\] где \(P_1\) --- окрестность \(a\) в \(\mathbb R^m\), \(Q\) --- окрестность \(b\) в \(\mathbb R^n\) («вместо шара --- произведение шаров, „полишар``; чтобы почувствовать природу этого множества, нужна хотя бы четырёхмерная картинка: круг в одной плоскости на круг в другой»).

\textbf{Шаг 3 (устройство обратного отображения).} Положим \[\Psi:=(\Phi|_W)^{-1}\colon V\to W.\] Так как \(\Phi\) \textbf{не меняет \(x\)-координату}, обратное отображение тоже её не меняет («точки, висящие над данным \(x\), не склеиваются, а деформируются эластично»): \[\Psi(u,v)=\bigl(u,\ h(u,v)\bigr)\qquad\text{для некоторого отображения } h\colon V\to\mathbb R^n\] (\(h\) --- вторая координата \(\Psi\)).

\textbf{Шаг 4 (множество \(P\)).} Ключевое наблюдение: точки, в которых \(F(x,y)=0\) --- это в точности точки, переходящие под действием \(\Phi\) на координатное подпространство \(\mathbb R^m\times\{0\}\): \[F(x,y)=0\quad\Longleftrightarrow\quad \Phi(x,y)=(x,0).\] Положим \[P:=\bigl\{u\in\mathbb R^m:\ (u,0)\in V\bigr\}.\] Так как \(\Phi(a,b)=(a,0)\in V\) и \(V\) открыто, точка \(a\) --- внутренняя для \(P\) (вокруг \((a,0)\) есть целый шарик из \(V\); пересечение открытого множества с подпространством открыто в подпространстве). \textbf{Требует сверки.} Формальная проверка открытости \(P\) именно в \(\mathbb R^m\) (через непрерывную проекцию/пересечение) --- та самая деталь, которую лектор отложил «на следующий раз» (см. ниже).

\textbf{Шаг 5 (определение \(\varphi\)).} Для \(x\in P\) положим \[\varphi(x):=h(x,0)\ =\ \text{вторая координата }\Psi(x,0).\] Тогда по построению \(\Psi(x,0)=\bigl(x,\varphi(x)\bigr)\), и, применяя \(\Phi\) (оно взаимно обратно \(\Psi\) на \(V\)): \[\Phi\bigl(x,\varphi(x)\bigr)=\bigl(x,\varphi(x)\bigr)\mapsto(x,0)\quad\text{точнее}\quad \Phi\bigl(x,\varphi(x)\bigr)=(x,0),\] а по определению \(\Phi\) вторая координата образа равна \(F\): \[F\bigl(x,\varphi(x)\bigr)=0\qquad\forall x\in P.\] «Если вы дочитали до этого места и не забыли, что обозначают буковки, то тут всё понятно, нет повода беспокоиться» (лектор). Кроме того \(\varphi(x)\in Q\): вторая координата точки \(\Psi(x,0)\in W=P_1\times Q\).

\textbf{Шаг 6 (гладкость и единственность).} \(\varphi\) --- композиция гладких отображений: \(x\mapsto(x,0)\) (вложение класса \(C^\infty\)) и \(\Psi=\bigl(\Phi|_W\bigr)^{-1}\) (класса \(C^r\) по теореме о гладкости обратного отображения, §9); значит \(\varphi\in C^r\). Единственность: \(\Phi|_W\) биективно, поэтому для данного \(x\in P\) среди точек \((x,y)\in W\) ровно одна переходит в \((x,0)\) --- а это в точности \((x,\varphi(x))\); значит, если \(y\in Q\) и \(F(x,y)=0\), то \((x,y)\in P_1\times Q=W\) переходит в \((x,0)\), и потому \(y=\varphi(x)\). \(\blacksquare_{(1),(2)}\)

\textbf{Картинка лектора.} Оси: «\(\mathbb R^m\)» (целое подпространство) и «\(\mathbb R^n\)». В \(W=P_1\times Q\) сидит точка \((a,b)\); диффеоморфизм \(\Phi\) пересаживает \(W\) в открытое \(V\subset\mathbb R^{m+n}\), причём \((a,b)\mapsto(a,0)\); слои \(\{x=\mathrm{const}\}\) деформируются эластично, но не склеиваются. Пересечение \(V\) с подпространством \(\mathbb R^m\times\{0\}\) --- открытое множество, содержащее \(a\); его прообраз при \(\Psi\) --- «поверхность», являющаяся графиком \(\varphi\): по каждому \(x\) игрек вычисляется однозначно.

\textbf{Требует сверки.} \textbf{Статус доказательства.} Лектор: «Я не проверил единственное {[}открытость \(P\) и финальные склейки{]}... В следующий раз проверю, заодно свежим взглядом. Это будет чисто формальная выкладка --- ещё раз из этих деталей я составлю какое-то равенство, в котором всё жёстко и однозначно. В некотором смысле это очевидная теорема {[}по модулю многомерности{]}... я отслеживаю самые простые вещи: проекция сюда, проекция туда, пересечение с подпространством --- ничего спекулятивного не делаю.» Итого: доказательство проведено по существу полностью, последняя формальная проверка отложена на следующую лекцию.
